% ============================================================
%  HCMUT-DEE Thesis Kit v2.0.4 — Résumé (French)
%  Copyright (c) 2026 Nguyễn Trọng Thắng & TS. Nguyễn Phúc Khải
%  License: MIT
% ============================================================

\chapter*{RÉSUMÉ DE THÈSE}
\addcontentsline{toc}{chapter}{Résumé de thèse}

La rédaction de rapports techniques et de thèses de fin d'études représente un défi majeur pour les étudiants en ingénierie en raison des exigences strictes en matière de mise en page académique. Les logiciels de traitement de texte traditionnels, comme Microsoft Word, révèlent souvent des limites importantes concernant la cohérence des styles, des bogues d'affichage multiplateformes et des difficultés à gérer le code source, les figures et les références bibliographiques. Pour résoudre ces problèmes, ce projet développe le **HCMUT-DEE Thesis Kit**, un modèle LaTeX optimisé conçu spécifiquement pour les étudiants de la Faculté d'Électricité et d'Électronique (FEEE) à l'Université de Technologie de Hô Chi Minh-Ville (HCMUT - VNU-HCM).

Ce document sert à la fois de guide complet d'utilisation et d'échantillon structurel officiel du modèle. Le kit est construit autour d'une classe de document robuste (`hcmut-dee.cls`) dérivée de la classe standard `report`, intégrant des paquets essentiels pour la mise en forme mathématique avancée, la génération automatique de la page de couverture (\texttt{\textbackslash makecover}), de la fiche de mission de thèse (\texttt{\textbackslash makeassignment}) et de la grille d'évaluation (\texttt{\textbackslash makeevaluation}). De plus, le modèle propose des styles de code préconfigurés pour Python, MATLAB, C/C++ et le langage PLC (Structured Text), répondant parfaitement aux besoins académiques des spécialités de l'électrotechnique, de l'électronique, des télécommunications et de l'automatique.

Les applications pratiques démontrent que le modèle réduit le temps de mise en page de plus de 80\%, élimine complètement les erreurs de marges débordantes ou de formatage des citations, et améliore considérablement le standard typographique et professionnel des rapports académiques. Ce projet incarne l'esprit de transmission et d'héritage des connaissances entre les générations d'étudiants et vise à promouvoir l'adoption de LaTeX dans la recherche scientifique et technique à l'université.

\vspace{0.5cm}
\noindent \textbf{Mots-clés:} Modèle LaTeX, Rapport technique, Thèse de fin d'études, Faculté d'Élec\-tri\-ci\-té et d'Élec\-tro\-nique, HCMUT, Héritage des connaissances.
